\documentclass{scrartcl}
\usepackage[utf8]{inputenc}
\usepackage{graphicx}
\usepackage{adjustbox}
\usepackage{amsmath}
\usepackage{float}
\usepackage{siunitx}
\sisetup{
  per-mode=symbol,
  separate-uncertainty
}
\usepackage{hyperref}
\hypersetup{
  colorlinks=true,
  linkcolor=blue,
  filecolor=magenta,
urlcolor=cyan,}
\usepackage{authblk}

\usepackage{minted}
\usepackage{xcolor}

\title{CGL-sheet12}

\author[1]{Robin Kühnberger}
\author[2]{Felix Naraschewski}

\affil[1]{University Stuttgart}
\affil[2]{University Stuttgart}

\date{\today}

\begin{document}
\maketitle
\tableofcontents

\section{Datentypen in C++}

\inputminted[breaklines, linenos, bgcolor=gray!10]{cpp}{datentypen.cpp}

\section{Zahlen Raten mit C++}

\begin{enumerate}
  \item \textbf{Präprozessor:} Verarbeitet Befehle mit \texttt{\#} (z.\,B. \texttt{\#include}), fügt Header ein und entfernt Kommentare. Das Ergebnis ist eine reine Übersetzungseinheit.

  \item \textbf{Compiler:} Übersetzt den bereinigten C++-Code in Assembler-Code und prüft dabei Syntax sowie Semantik.

  \item \textbf{Assembler:} Wandelt den Assembler-Code in Maschinencode (Objektdateien, \texttt{.o}) um, der vom Prozessor lesbar, aber noch nicht ausführbar ist.

  \item \textbf{Linker:} Verbindet die Objektdateien mit nötigen Bibliotheken (z.\,B. für \texttt{std::cout}) zur finalen, ausführbaren Datei.
\end{enumerate}

\begin{minted}[breaklines, bgcolor=lightgray!10]{bash}
g++ -fdiagnostics-color=always -g -Wconversion -std=c++20 -O3 /path/to/main.cpp -o /path/to/main
\end{minted}

\begin{itemize}
  \item \textbf{Build:} Übersetzt \texttt{main.cpp} mit modernem C++20 in die ausführbare Datei \texttt{main}.
  \item \textbf{Einstellungen:} Kombiniert maximale Geschwindigkeitsoptimierung (\texttt{-O3}) mit Debug-Informationen (\texttt{-g}).
  \item \textbf{Sicherheit:} Aktiviert farbige Ausgaben und warnt vor potenziellen Datenverlusten bei Typumwandlungen.
\end{itemize}

\inputminted[breaklines, linenos, bgcolor=gray!10]{cpp}{zahlen-raten.cpp}

\section{Funktionen in C++}

\begin{minted}[breaklines, bgcolor=gray!10]{cpp}
bool pruefeTipp(int tipp, int loesung)
\end{minted}

\begin{itemize}
  \item \textbf{Rückgabetyp:} Der Datentyp ist \texttt{bool} (wahr/falsch).
  \item \textbf{Anzahl Parameter:} Es werden \textbf{2} Parameter übergeben (\texttt{tipp} und \texttt{loesung}).
\end{itemize}

\inputminted[breaklines, linenos, bgcolor=gray!10]{cpp}{funktionen.cpp}

%
%\begin{figure}[H]
%  \centering
%  \includegraphics[width=1\textwidth]{fibonacci.pdf}
%  \caption{Golden ratio}
%  \label{img:04}
%\end{figure}

%\section*{Quellen}
%\bibliographystyle{plain}
%\bibliography{bibliography}

%\listoffigures
%\listoftables

\end{document}